\chapter*{Preface}

These notes correspond to lectures given
\begin{itemize}
\item at \emph{Institut d'Astrophysique de Paris} (France) in March-April 2016, within the
framework of the \emph{IAP Advanced Lectures}:\\
{\footnotesize\url{https://www.iap.fr/vie_scientifique/cours/cours.php?nom=cours_iap&annee=2016}}
\item at the \emph{Centre for Cosmology, Particle Physics and Phenomenology} in Louvain-la-Neuve
(Belgium) in November 2016, within the framework of the \emph{Chaire Georges Lemaître}:\\
{\footnotesize \url{https://uclouvain.be/fr/instituts-recherche/irmp/chaire-georges-lemaitre-2016.html}}
\item at the
\emph{Bogoliubov Laboratory of Theoretical Physics}, in Dubna (Russia) in May 2017,
within the framework of the \emph{Dubna International Advanced School of Theoretical Physics}:\\
{\footnotesize\url{http://www.jinr.ru/posts/lecture-course-geometry-and-physics-of-black-holes/}}
\item at the Summer School \emph{Gravitational Waves 2018}, taking place at Les Houches
(France) in July 2018:\\
{\footnotesize\url{https://www.lkb.upmc.fr/gravitationalwaves2018/}}
\item remotely at the \emph{School on Black Holes and Gravitational Waves}
organized at the
\emph{Centre for Strings, Gravitation and Cosmology} of the
\emph{Indian Institute of Technology Madras}, Chennai (India) in January 2022:\\
{\footnotesize\url{https://physics.iitm.ac.in/~csgc/events/sbhgw}}
\item at \emph{École Normale Supérieure}, Paris (France) in May-June 2023, as
part of the PSL graduate programs in Physics and in Astrophysics:\\
{\footnotesize\url{https://relativite.obspm.fr/blackholes/paris23}}
\end{itemize}

\vspace{2ex}

In complement to these notes, one may recommend various monographs
devoted to black holes:  O'Neill (1995) \cite{ONeil95}, Poisson (2004) \cite{Poiss04}, Frolov \& Novikov (1998) \cite{FroloN98},
Frolov \& Zelnikov (2011) \cite{FroloZ11}, Bambi (2017) \cite{Bambi17} and Chru\'sciel (2020) \cite{Chrus20}, as well as the review
articles by Carter (1987) \cite{Carte87}, Chru\'sciel (2002, 2005) \cite{Chrus02, Chrus05},
Chru\'sciel, Lopes Costa \& Heusler (2012) \cite{ChrusLH12} and Wald (2001) \cite{Wald01}.
In addition, let us point other lecture notes on black holes:
Andersson, Bäckdahl \& Blue (2016) \cite{AnderBB16}
Compère (2006, 2019) \cite{Compe06,Compe19},
Dafermos and Rodnianski (2008) \cite{DaferR13},  Deruelle (2009) \cite{Derue09}
Hawking (1994) \cite{Hawki94,HawkiP15}, Reall (2020) \cite{Reall20} and Townsend (1997) \cite{Towns97}.

The history of black holes in theoretical physics and astrophysics is
very rich and fascinating. It is however not discussed here, except in some
small historical notes. The interested
reader is referred to Nathalie Deruelle's lectures \cite{Derue09}, to Kip Thorne's
book \cite{Thorn94}, to Carter's article \cite{Carte06}
and to Jean Eisenstaedt's articles \cite{Eisen82,Eisen93}.


The web pages associated to these notes are
\begin{center}
\url{https://relativite.obspm.fr/blackholes}
\end{center}
They contain supplementary material, such as the SageMath notebooks presented in
Appendix~\ref{s:sam}.

\vspace{2ex}

I warmly thank Cyril Pitrou for the organization of the Paris 2016 lectures,
Fabio Maltoni and Christophe Ringeval for the organization of the Louvain-la-Neuve ones,
Anastasia Golubtsova and Irina Pirozhenko for the
organization of the Dubna ones, Bruce Allen, Marie-Anne Bizouard, Nelson
Christensen and Pierre-François Cohadon for the organization of Les Houches ones,
Chandra Kant Mishra for the organization of the Chennai ones,
as well as Jean-François Allemand for the organization of the Paris 2023 ones.

\vspace{3ex}
These notes are released under the
\begin{center}
\href{https://creativecommons.org/licenses/by-nc-sa/4.0/}{{Creative Commons Attribution-NonCommercial-ShareAlike 4.0 International License}}\\[1ex]
\includegraphics[height=0.03\textheight]{cc_license.png}
\end{center}

